%%-------------Abstract-----------------
% Length:Ideally, one page or less
% Ph.D. Dissertations: 350 word maximum.
\abstract{ Microbes can transform their environment in unpredictable ways with beneficial outcomes. Microbial species in nature are rarely found in isolation. Instead most of them form complex communities or consortia with other microbial species. The transformative potential of such microbial communities far exceeds that of any individual species. Microbial co-cultures grown under controlled laboratory environments or microbial consortia found in nature, by virtue of their own biology, can produce molecules with applications in multiple domains e.g. medicine, agriculture and bio-energy. Hence, a structured approach to the understanding of synergistic microbial capacity can significantly boost the utilization of microbial co-cultures/consortia for industrial applications.
% Microbes have significant ability to transform their environments, but the metabolic capacity of a single microbial species is usually limited. However, microorganisms do not naturally exist in isolation -- they form co-cultures with other microbes that together exhibit complex, often unpredictable, metabolic capacities. Co-culture metabolisms have tremendous potential for many medical, environmental, and industrial applications, but harnessing them requires a systems-level understanding of synergistic microbial capacity.
To this end, a modelling framework has been developed that can demonstrate the difference in metabolism between the individual mono-cultures and the co-culture of two simplistic {\it{in silico}} organisms. This framework can produce hybrid automata representations of an organism directly from mass balance, stoichiometry, and flux bounds. This approach is both generic with respect to the microbe(s) under consideration and scalable in the size of the relevant metabolic network(s). Constructing a reduction from the trajectory computed by dynamic flux balance analysis (dFBA) to a hybrid automaton with the same dynamical behavior comprises the core of the modelling effort. 
% dFBA is a well-studied algorithm in systems biology to simulate microbial growth in batch cultures. It uses an iterative approach based on a quasi-steady-state assumption of concentration rates to evolve the organism across a discrete time horizon. 
The transformation to hybrid automaton is achieved by basis vector decomposition of the flux cone and the subsequent use of non-determinism to generically represent trajectories. Considerations from control theory are also taken into account. Using Ariadne, a C++ library for hybrid systems analysis, two in silico organisms are evolved in both mono-culture and co-culture, identifying emergent metabolic phenomena resembling the predictions given by dFBA. Furthermore, this framework provides a foundation to perform parametric analysis and identify unsafe regions.
% Dynamic flux balance analysis is a well-established algorithm in systems biology that combines stoichiometry-based FBA with the dynamics of the growth media  to compute the metabolism of an organism in a batch culture. 
%In this effort, we develop a reduction from dynamic flux balance analysis to a hybrid automaton that exhibits the same trajectories in the metabolic flux space. The reduction to hybrid automaton is achieved by decomposing the steady state metabolism, where metabolic fluxes are decomposed into basis vectors, along with the introduction of clock variables to determine the upper and lower bounds of the flux at each state. Our approach was demonstrated using a hybrid systems tool Ariadne and yielded results that show satisfaction of reachable production of emergent metabolites.  Additional functionalities that are enabled are visualization of the evolution of the metabolites, identification of unsafe regions and exhaustive traceability analysis.
}